\emph{Carried over verbatim from the full report (draft v2.4, 23 September 2026), which is released with this report. Section and appendix numbers in this appendix refer to that report; Appendix~E maps them to this one.}

Validation that fits on some problems and tests on others (H1, H10, H11b, H20, H23, H36) is grouped by problem: no problem appears on both sides of a split. The cross-model tests hold out the model but not the problem. In H7, 234 of the 335 test rows have their problem in training, 198 of them with the same label and base accuracy, including all 52 DeepSeek rows, which are exact copies of training rows published as \texttt{qwen3-8b:low}, so part of that test is in-sample (with them counted once, H7 is still +0.0). In H12 to H15, 70 of the 109 test pairs have their problem in training, where the classifier is trained on the same DeepSeek feature vectors labeled by other models, an overlap whose effect was not measured and need not favour the hypotheses, since another model's label on the same problem did worse than the per-model majority in H3. H7, H12, H14 and H15 died anyway, and H13 is withdrawn. Warmup rows are as in v1.0 except where later drafts corrected them (Appendix D, items 6, 21, 24 to 29, 51 and 54). Draft v2.2 made H2 die once the duplicated rows are counted once, moved H13 from dead to a withdrawn oracle ceiling, qualified H16 as within re-execution noise, moved H17 from survives to not established, marked H6 and T as having no kill criterion and corrected the descriptions of H3, H17 and T; this draft describes the classifier actually run for H10, corrects the description of H14, marks H18 as having no kill criterion, adds to H2 the result of E1's original test and to H13 its value without constant predictions, gives H1, H7 and T with the duplicated rows counted once, gives H7 against the baseline its pre-registration named, and describes the rows on which H17's threshold was fitted. Main-phase hypotheses are new.

{\footnotesize\setlength{\tabcolsep}{4pt}
\begin{longtable}{L{1.1cm}L{3.7cm}L{5.7cm}L{2.2cm}}
\toprule
\textbf{ID} & \textbf{Hypothesis} & \textbf{Result versus its baseline} & \textbf{Verdict} \\
\midrule
\endfirsthead
\toprule
\textbf{ID} & \textbf{Hypothesis} & \textbf{Result versus its baseline} & \textbf{Verdict} \\
\midrule
\endhead
\midrule
\endfoot
\bottomrule
\endlastfoot
H1 & Problem text predicts robustness (32 features, grouped 5-fold) & 0.675 vs 0.761, $-$8.7 (0.682 vs 0.756, $-$7.4, with the duplicated rows counted once) & dead \\
H2 & Robustness is a property of the problem (cross-model concordance) & z = 0.88 and $-$1.3 points against the per-model majority (bar 4.0) once the 52 rows published under both \texttt{qwen3-8b:low} and the DeepSeek name, identical in label and base accuracy, are counted once, and z = 1.81 under E1's original global-rate test (bar 2); z = 4.04 and +3.4 only when they are counted twice & dead \\
H3 & Model A's label predicts model B (56 ordered model pairs over 10 problems, 560 predictions) & $-$7.3 & dead \\
H4 & A's base accuracy predicts B's robustness & $-$11.8 & dead \\
H5 & Difficulty order transfers (half-batch rule) & $-$25.4 & dead \\
H6 & Same, oracle k & $-$13.0 & ceiling, negative (no kill criterion) \\
H7 & Base accuracy beats the prior on held-out models & +0.0 against the global prior of the training models, used where PREREG\_E4 named the per-model majority ($-$1.2 against that; +0.0 and $-$1.4 with the duplicated rows counted once) & dead \\
T & Trivial-robust rule (robust iff base acc < 0.05), formulated after the data of H9 & +6.0 over the 335 published pairs (+7.1 with the duplicated rows counted once), +0.0 on \texttt{val-sample} & trap (no kill criterion) \\
H10 & A classifier on the official uncertainty features reaches the accuracy the organizers' proposal reports for their baseline (grouped CV) & 0.7744 against a majority constant of 0.756, inside the 0.692 $\pm$ 0.101 that the proposal reports (best of several classifiers, stratified 10-fold CV on its sample validation set); run as a logistic regression on 13 of the 14 features, three of them constant in our extraction, where PREREG\_E6 specified the official regressor and threshold & confirmed (method substituted; not declared at the time) \\
H11b & Length, saturation and trend features add signal & $-$1.8 / $-$2.5 & dead \\
H12 & Cached model's real uncertainty features predict other models & $-$16.5 & dead \\
H13 & Same, with a threshold chosen per target on its own items & +7.3 (+5.5 when constant predictions are excluded), all from thresholds chosen on the items they are scored on; $-$2.8 once the number of predicted positives is fixed at the true count (E8) & ceiling, withdrawn (no kill criterion) \\
H14, H15 & Calibration by the target's public robust rate (H14) or by the global rate (H15) & $-$2.8, $-$10.1; each target's public rate is its rate on the tested items, so H14 coincides with the oracle-k ceiling of E8 & dead \\
H16 & Self-measurement of every model beats the constant & 10/17 > 9/17 on one execution, one item, inside the spread of H18 & survives its criterion, within noise \\
H17 & A gpt-oss threshold calibrated on the ten gpt-oss rows of \texttt{val-sample}, one of them robust (6.2), other models unchanged, beats H16 & 11/17 > 10/17 on its first execution; three re-executions of the same bundle gave 9, 11 and 11, and 9/17 meets the kill criterion & not established \\
H18 & Identical bundles score differently across executions & \{11, 9, 11, 11\} & confirmed (no kill criterion: PREREG\_V7 fixes the measurement, not a criterion) \\
H19 & Our five mechanical perturbation families reproduce the official label (DeepSeek, 99 pairs) & accuracy 0.354 vs 0.788, kappa \textbf{0.02} & dead \\
H20 & A threshold tuned on our label transfers to theirs & 0.212 vs 0.788 & dead \\
H22 & Translation-stability (it/fr/de) reproduces the official label (77 pairs) & 0.507 vs 0.740, kappa \textbf{$-$0.10} & dead \\
H23 & A threshold tuned on translation-stability transfers & +0.0 & dead \\
H29 & Representation-distance features are repeatable across batch size and padding (gate; E11 plan, not a frozen pre-registration) & ICC(2,1) 0.997 to 0.999, from the dated outcome file of E11 (the per-item output was written to the Modal volume and is not in the local record) & passes \\
H30 & Residualized representation shift under perturbation predicts the label (E11 plan) & AUROC 0.615: below the death threshold of 0.655 fixed in the plan (the null p95 over its planned cells), above 0.604, the uncorrected p95 for a single statistic (permutation p = 0.030); n\_char 0.705 & dead \\
H31 & Directional coherence of the shift predicts the label (E11 plan) & 0.555 (permutation p = 0.29) & dead \\
H33 & Self-consistency over 8 samples predicts the label (164 pairs) & AUROC \textbf{0.502} & dead \\
H32 & The organizers' families, regenerated, reproduce their label (157 pairs; four of five families, \texttt{domain} excluded by the pre-registered 60\% acceptance gate) & kappa \textbf{$-$0.001} ($-$0.030 with all five) & dead \\
H35 & The organizers' published probe artifact orders our ground truth & 0.709 with its own inference aggregation, mostly problem length (0.57 once residualized on log character count); 0.739 with the pre-registered one (0.62 once residualized, as for our probe); its control probes, trained on shuffled labels, give 0.68 by themselves at layer 31 (8.3) & survives, largely explained \\
H36 & Our linear probe, nested CV, orders our ground truth & 0.710 OOF against a reduced nested null (p95 0.574 with 60 permutations and two of four regularization values, where 500 and the full grid were pre-registered; 0.570 with the full grid when re-run on 23 September); about +0.02 over character count in a second-level cross-validation (8.3) & survives, largely explained \\
E13b-1 & Under the main-phase rule our regenerated label is test-retest reliable (kappa(A, B) $\geq$ 0.60 on at least 3 of 4 models) & 0.34 to 0.56, none reaches 0.60 (7.5) & dead \\
E13b-2 & Under the main-phase rule it agrees with the official label (DeepSeek, kappa $\geq$ 0.40) & 0.02 [$-$0.09, 0.13] (7.5) & dead \\
E13b-3 & The labeling rule drives the instability (kappa(A, B) gain $\geq$ 0.15 on at least 3 of 4 models) & gain $\geq$ 0.15 on 2 of 4, and on 0 of 4 when both rules are read on the same problems (7.5) & dead \\
\end{longtable}
}

Of the thirty-one entries in the table (thirty rows; H14 and H15 share one), twenty-seven carried a kill criterion written before they ran: twenty-one died by it, five survived it (H10, H16, H29, H35, H36) and one (H17) is not established. Three of the twenty-seven (H29 to H31) had their criterion in the written plan of E11, fixed before the runs, rather than in a frozen pre-registration. H2 is weaker than the rest: its concordance (56 of 73) had already been observed when PREREG\_E1 was written, and the per-model null and the bar used in its row were added in E3, after the first run; with the duplicates counted once, H2 also fails E1's original test (z = 1.81 against 2). The other four entries have no kill criterion of their own: H6 and H13 are oracle ceilings (H6 came out negative; H13 came out +7.3 and was withdrawn in E8), T is a rule formulated after seeing the data of H9, stopped by the check on \texttt{val-sample} announced in E4, and H18 is a repeat-run measurement whose pre-registration fixes what to measure but no criterion.

Executed but not in the table: H8, the E4 oracle ceiling (0.776, +1.5 points against the global constant and +0.3 against the per-model majority; +1.8 and +0.4 with the duplicated rows counted once); H9, whose outcome became rule T; and the V4 allocation entry (909843, Appendix B), whose pre-registration stated that a tie at 9/17 was impossible unless the model path on the DeepSeek items failed. It tied at 9/17. None of the three explanations recorded then was verified; the first, drift on the DeepSeek items, is the kind of re-execution noise later measured in H18. H1, H7, H8 and T were computed on all 335 published pairs, so the 52 duplicated DeepSeek rows count twice; counted once (283 pairs), no verdict changes, and the values are given above. Predictions attached to measurement experiments are reported where they belong: against our written predictions, batch size 1 does not stabilize generation and the configuration without compilation lands in a state of its own (8.1). E15, the experiment behind H35 and H36, also carried a written prediction, that both probes would die; both survived their null (8.3).

E11 departed from its plan in ways not declared at the time. The plan listed three further primary tests, noise injection (H32), a geometric control (H33) and a trajectory-angle control (H34), which were never run; distributional stability on the unlabeled models (H35) was conditional on a channel passing and was correctly not run. The plan required Holm over five primary tests at alpha 0.010; the analysis script announces Holm over the two tests it ran but computes no correction, and H30, with a permutation p of 0.030, would not have passed Holm over those two even at the alpha of 0.05 implied by the script's p95 threshold (0.030 against 0.025). The analysis used all 164 problems, where the plan held out 40 as a confirmation set and ran the primary analysis on 124. It computed each channel's AUROC directly on the scalar, with the orientation chosen on the data, rather than out of fold from an L2 logistic regression with C = 0.05 over 5 folds and 20 repetitions. It replaced the six-feature text baseline of H28 with the AUROC of character count alone. It omitted the two analyses that the plan (its section 4.1) made obligatory next to every reported AUROC, the same AUROC stratified on quartiles of the number of perturbed predictions and the reference AUROC of the bookkeeping counters (0.862 out of fold, not available at inference). And it extracted only the primary depth (0.80) of the three the plan fixed. On the 164 problems analyzed, H30 is below the plan's threshold of 0.655 and trails character count by 0.09, and H31 does not reach its null. Identifiers H32, H33 and H35 were reused in E12 and E15 for the hypotheses listed here; H34 was not reused. H21 and H24, the second stages of E9 and E10, were to run only if H19 and H20, and H22 and H23, survived, and were correctly not run. The plan's H25 to H28 are not in the table: H25 restated as a confirmation the A100 anomaly that 8.2 still reports (AUROC 0.36 in the rule's orientation), traced to the feature that our own extraction returned as 1.0 on every item (2.3b), and its prospective part, a simulation with a kill criterion, has no recorded outcome; H26 (the aggregation defect of 2.3a) and H27 (a label-noise ceiling) were measured before the plan was written; and H28, a text-only baseline, was also measured before the plan (character count 0.687 out of fold, as recorded there), while its planned non-inferiority test against the official features was never run.

Two of the three methodological notes of v1.0 stand: in-sample signal refused to generalize twice, and there is the distribution trap of rule T (2.4). The third, that grouped validation is what killed H1, is withdrawn (Appendix D, item 12): under an ungrouped 5-fold split, with the same code and data, the same text model also stays below the majority constant ($-$0.9, $-$0.3 and $-$2.1 points over three seeds, against a bar of 2.7), so grouping deepened the loss but did not cause it. The main phase added a fourth, in H30: a control matched on the nuisance variable and decided before the data (token-matched neutral padding) separated a length signal from a fragility signal; in 8.3 the same separation needed a control we added only after seeing the data, which is why such a control belongs in the pre-registration. The verification of draft v2.1 added a fifth: the public datasets publish one measurement under two model names, and counting it twice produced the only surviving warmup hypothesis about a property of the problem (H2).

Sections 4.1 to 4.3 of draft v1.0 are not reproduced in this draft; their content is in the dated files PREREG\_E1 to PREREG\_E8 and PREREG\_E5\_ESITO, which will be released with the final version (Section 13).
